\section{Scope and limits}\label{sec:limits}

The preceding sections established a chain: a growth law for the denominators,
a local classification, a moduli space, a located point, the removal of a
projective ambiguity, a point of the character variety, and transport along the
flow. This section states what that chain does and does not establish. It is
written so that a reader who stops here is not left to infer the strength of
any link from the confidence of the prose.

\subsection{Membership: a caution, and its discharge}\label{ssec:verdict}

The membership result was obtained twice, and the order in which it was
obtained is part of what it means.

\begin{remark}[Membership, with the original caution; \texttt{S16-14}, \textbf{STRUCTURAL}]
System $\mathsf C$ satisfies every hypothesis under which van der Put and Saito
build the moduli space $\mathcal M$ for the $(0,0,1)$ row, it is a point of the
chart $U_0$, and since Painlev\'e V \emph{is} the isomonodromy foliation on
that space, and since $q(q-1)\neq0$ and $t\neq0$ at the point so that the
vector field is regular there, exactly one leaf passes through System
$\mathsf C$. The deformation exists and is not an ansatz. The caution attached
to it was this: the match was of \emph{projective} normal forms, so what was
established was that $\mathsf C$'s projectivisation is that of the family
member at $(p,q,t)$ --- equality up to a rank-one twist, in the same sense in
which $\mathsf B$ and $\mathsf C$ are projectively equivalent. That caution is
the reason the label is \textbf{STRUCTURAL} and not \textbf{PROVEN}.
\end{remark}

The caution was real, and it was not overlooked: it was stated by the session
that produced the result, against its own conclusion. It was discharged
afterwards, by a different computation.

\begin{remark}[The caution removed; \texttt{S17-04}, \textbf{PROVEN}]
A second-order operator is determined by the \emph{pair} $(P,Q)$ in
$\partial^2+P\partial+Q$. The projective invariant
$R=P^2/4+P'/2-Q$ that the earlier match used is blind to a rank-one twist,
whereas $(P,Q)$ is not. Comparing both coefficients: System $\mathsf C$ in the
$z$ coordinate has
\[
  P_{\mathsf C}=\frac{2z-1}{z(z-1)},\qquad
  Q_{\mathsf C}=\frac{22z^{2}-22z-2i\sqrt{11}\,z+5+i\sqrt{11}}{54\,z(z-1)},
\]
and the family's scalar operator at the located point $(p,q,t,\theta_\infty)$
with cyclic vector $e_1$ has exactly the same $P$ and $Q$:
$P_{\mathsf F}-P_{\mathsf C}=0$ and $Q_{\mathsf F}-Q_{\mathsf C}=0$
identically. The two differential modules are therefore \emph{isomorphic}, not
merely projectively isomorphic.
\end{remark}

\begin{remark}[The consequence, and its ground; \texttt{S17-06}, \textbf{STRUCTURAL}]
The ``up to a rank-one twist'' caveat is accordingly withdrawn. It is withdrawn
because the full operators were verified to agree, and \emph{not} because
trace-freeness singles out a representative of a projective class; that
argument is false, and \S\ref{ssec:gauge} says why. The defensible wording is
that the differential module of System $\mathsf C$ is the module of the van der
Put--Saito family at the located point of $U_0$, with no twist qualifier.
\end{remark}

Two things did not change when the caveat was withdrawn.

\begin{remark}[The statement changed, not the label; \texttt{S17-08}, \textbf{STRUCTURAL}]
Membership was deliberately not raised above \textbf{STRUCTURAL}. Any such
raise was tied in advance to a falsifier that remains unrun
(\S\ref{ssec:unrun}). What the discharge changed is the \emph{statement} ---
membership is now about System $\mathsf C$'s actual differential module rather
than its projectivisation --- and not the strength of the evidence for it.
\end{remark}

This is the reason the paper's central result is labelled as it is. The
operator identity is \textbf{PROVEN}; the membership statement built on it is
\textbf{STRUCTURAL}. Nothing in the intervening sections upgrades the second
because the first was strengthened.

\subsection{The falsifier that was not run}\label{ssec:unrun}

\begin{remark}[Not established; \texttt{S16-15}, \textbf{CONJECTURED}]
The pre-registered falsifier was \emph{not run}. It requires integrating
Painlev\'e V from $(t_0,q,p)$ to $t_0+\varepsilon$ for two values of
$\varepsilon$, rebuilding the linear system, and computing monodromy
\emph{and Stokes} data agreeing with $\mathsf C$'s to at least $30$ digits.
Computing Stokes data for an irregular rank-two system at that precision is a
separate undertaking, and it was outside the budget. What was done instead, and
what is reported as supplementary rather than as the falsifier, is a regularity
check of the vector field at the point.
\end{remark}

The threshold was left unchanged for whenever the falsifier is run. No
threshold was weakened to close this gap, and the gap is recorded as a gap. A
reader should take the label at face value: \textbf{CONJECTURED} here means
that a specific, stated, still-available test has not been performed.

\subsection{What the numerical confirmation covers}\label{ssec:numerical-scope}

The transport computation of \S\ref{sec:leaf} closed the same gap
\emph{numerically}, which is a different thing, and the difference is the
subject of this subsection.

\begin{remark}[Constancy, and its limits; \texttt{S22-15}, \textbf{VERIFIED}]
The trace coordinates $x_1,x_2,x_3$ are constant along the flow to $90.81$
digits at $\varepsilon=10^{-3}$ and $91.83$ digits at $\varepsilon=10^{-4}$,
against a threshold of $30$ digits. This is a numerical agreement at one point
along one direction in $t$; \emph{it is not a proof of leaf membership}.
\end{remark}

That last sentence belongs to the claim and not to the drafting. Three further
qualifications delimit it.

First, the confirmation is for the trace coordinates, and for those alone. It
is not a confirmation of the monodromy representation itself, still less of the
Stokes data, which is what \S\ref{ssec:unrun}'s falsifier would require.

Second, it is a confirmation \emph{under one convention}, and the selection of
that convention is itself a result rather than an assumption.

\begin{remark}[The convention is a result; \texttt{S22-11}, \textbf{VERIFIED}]
Of $16$ candidate conventions exactly one satisfies all four pre-registered
conditions --- right action, small loops counter-clockwise, $M_\infty$
clockwise, van der Put--Saito's $M_1=\gamma_1$. For it,
$\lVert M_0M_1M_\infty-I\rVert=2.44\times10^{-105}$; the other fifteen fail,
several on $\alpha$ alone, at $0.327$, $2.27$ and $2.34$.
\end{remark}

A reader is entitled to ask whether the result is an artefact of a convention
chosen to make it work. The answer is that the convention was fixed by
pre-registered conditions before the comparison, that fifteen of sixteen
alternatives are excluded by those conditions, and that the exclusion is
quantitative rather than stylistic. What is \emph{not} shown is that the result
would survive a different set of conditions; that is the open question of
\S\ref{ssec:convention-independence}.

Stated for the record, in the form in which it should be quoted:
\emph{membership is confirmed for the trace coordinates $x_1,x_2,x_3$, under
the one convention of sixteen that passed, and the selection of that convention
is itself a result and not an assumption.}

Third, the agreement does not improve without bound, and the reason is
understood.

\begin{remark}[The residual is systematic, not round-off; \texttt{S22-17}, \textbf{VERIFIED}]
Raising the working precision from $200$ to $320$ digits leaves the agreement
and its ratio unchanged to every printed figure --- $90.8118$, $91.8301$, ratio
$10.429999$ --- while every precision-limited control improves by about $120$
digits. Over $\varepsilon=10^{-3},10^{-4},10^{-5}$ the ratios are $10.430$ and
$10.034$. The residual is therefore a systematic $O(\varepsilon)$ error of the
integrator, not round-off, and it is floored by the optimal-truncation level
$10^{-88.7}$ set by the pre-registered truncation radius, which no increase in
precision can reduce.
\end{remark}

The practical consequence is that ninety digits is the instrument's ceiling at
this $\varepsilon$, and that quoting a larger figure would require a different
integrator rather than a longer run.

The second branch of the sign pair was treated the same way rather than
inferred from the first.

\begin{remark}[Branch two was run; \texttt{S22-18}, \textbf{VERIFIED}]
Branch $2$, at $t=+2\sqrt{33}\,i/9$, passes the convention selection, the
sensitivity stage and the transport stage independently: exactly one convention
of $16$ is selected, the sensitivity ratios are $99.9952$ and $99.9998$, and
the transport agreement is $90.82$ and $91.84$ digits against the $30$-digit
threshold. It was run, not inferred.
\end{remark}

\subsection{What this paper does not claim}\label{ssec:not-claimed}

Collected, so that they are not left to be reconstructed:

\begin{enumerate}
\item That leaf membership is \emph{proven}. It is \textbf{STRUCTURAL}, and the
  test that would raise it is named and unrun (\S\ref{ssec:unrun}).
\item That the numerical transport constitutes such a proof. Its own claim says
  it does not.
\item That the monodromy or Stokes data of System $\mathsf C$ have been
  compared with the family's along the flow. They have not.
\item That the projective normalisation has a unique representative on general
  grounds. The natural argument for that is false
  (\S\ref{ssec:gauge}).
\item That the result is independent of the monodromy convention. One
  convention was selected by stated criteria; independence is open
  (\S\ref{ssec:convention-independence}).
\item That any displayed formula in this paper has been re-derived from the
  computation that produced it. Each traces to a recorded claim; none has been
  symbolically re-derived here, and the formulas of
  Appendix~\ref{app:typo} were transcribed from page images rather than
  computed.
\end{enumerate}
